\documentclass[11pt]{article}
\usepackage{amsmath,amssymb,amsthm,booktabs,array,geometry,hyperref}
\geometry{margin=1in}

\newtheorem{theorem}{Theorem}
\newtheorem{proposition}[theorem]{Proposition}
\newtheorem{definition}[theorem]{Definition}
\newtheorem{remark}[theorem]{Remark}

\title{LEAd Architectural Decision:\\
Second-Argument Convention, F16$^+$ Policy, and Final Recommendation}
\author{Monogate Research}
\date{2026-04-21}

\begin{document}
\maketitle

\begin{abstract}
LSE\_Lower\_Bound.tex identified an open architectural question:
should $\mathrm{LEAd}(a,b) = \ln(e^a + b)$ (with raw second argument $b$)
be promoted to a genuine F16 primitive, creating an extended F16$^+$ model
where $N_+(\mathrm{LSE}) = 2$ instead of 4?
This paper delivers the formal architectural decision.
We weigh five criteria: (1) group-theoretic consistency with G$_{16}$,
(2) impact on SuperBEST core and LSE, (3) library implementation benefit,
(4) theorem-statement coherence, and (5) precedent in the 23-op taxonomy.
\textbf{Decision}: LEAd is \emph{not} promoted to the standard F16 primitive set.
It is recognized as a genuine computational shortcut and assigned to Layer 2
as the 24th operator — but designated as a \emph{structural extension}, not an
F16 orbit element. This creates a clean separation: 16 orbit operators (F16),
7 algebraic-shortcut/notation operators (the current extended set), and 1
structural-extension operator (LEAd). Theorem statements use F16-24 = F16 for
Layer 1; LEAd is available for Layer 2 library use.
A CONTEXT.md update is included.
\end{abstract}

\section{The LEAd Operator: Recap}

\begin{definition}[LEAd]
$\mathrm{LEAd}(a, b) = \ln(e^a + b)$, where $b > 0$ is the ``raw'' second argument:
it enters the computation directly (not via $\ln(b)$).
\end{definition}

\begin{definition}[Standard F16 argument convention]
In the standard F16 orbit, every operator has:
\begin{itemize}
  \item First argument $a$: via $\exp(\pm a)$
  \item Second argument $b$: via $\ln(\pm b)$ or $1/\ln(b)$
\end{itemize}
LEAd violates this: its second argument $b$ appears raw (not via $\ln$).
LEAd is therefore outside the G$_{16}$ orbit of EML.
\end{definition}

\begin{proposition}[Key impact: LSE]
\begin{itemize}
  \item Without LEAd (standard F16): $N(\mathrm{LSE}) = 4$
    (exp$+$exp$+$add$+$ln, proved tight in LSE\_Lower\_Bound.tex)
  \item With LEAd (F16$^+$): $N_+(\mathrm{LSE}) = 2$
    (exp$(y)$[1n] $+$ LEAd$(x, e^y)$[1n], proved numerically)
\end{itemize}
Impact: 2n saving on LSE in any F16$^+$ model.
Impact on downstream: partition function pair drops from 10n (F16) to 3n (F16$^+$)
if we chain LEAd with exp.
\end{proposition}

\section{Five-Criterion Analysis}

\subsection{Criterion 1: Group-Theoretic Consistency}

The F16 orbit is generated by $G_{16} \cong (\mathbb{Z}/2)^4$ acting on
EML$(x,y) = e^x - \ln y$. The four generators are:
negate first arg, invert second arg (replacing $\ln(b) \to -\ln(b) = \ln(1/b)$),
negate result, take reciprocal of result.

\textbf{LEAd does not arise from this action}: no combination of these generators
transforms the second-arg slot from $\ln(b)$ to $b$ (raw).
The $G_{16}$ action preserves the argument structure; LEAd requires changing
the argument type itself.

\textbf{Verdict}: promoting LEAd to F16 breaks $G_{16}$ symmetry.
The resulting ``F16$^+$'' orbit has 17 elements but is no longer the orbit of a
single operator under a group action — it is an ad hoc extension.

\subsection{Criterion 2: Impact on SuperBEST Core}

\begin{center}
\renewcommand{\arraystretch}{1.2}
\begin{tabular}{lccc}
\toprule
Expression & F16 (no LEAd) & F16$^+$ (with LEAd) & Change \\
\midrule
SuperBEST core 15n/79.5\% & \textbf{unchanged} & \textbf{unchanged} & None \\
LSE$(x,y)$ & 4n & 2n & $-2n$ \\
Partition fn.\ pair & 10n & 3n & $-7n$ \\
Softplus & 2n & 2n & None \\
Von Neumann entropy & 7n & 5n & $-2n$ \\
Boltzmann ratio & 6n & 6n & None \\
\bottomrule
\end{tabular}
\end{center}

The 15n/79.5\% core result is \emph{unaffected} by LEAd (it uses arithmetic operators,
not LEAd). The impact is concentrated in expressions involving $\ln(e^x + e^y)$ form.

\subsection{Criterion 3: Library Implementation Benefit}

$\mathrm{LEAd}(a, b) = \ln(e^a + b)$: this is a genuine fused operation with
real numerical benefit:
\begin{itemize}
  \item \textbf{Numerical stability}: avoids computing $e^a$ and $e^y$ separately then adding,
    which can overflow for large $a$. A fused implementation can use the log-sum-exp trick
    internally ($\ln(e^a + b) = a + \ln(1 + b e^{-a})$ when $a$ is large).
  \item \textbf{Performance}: one fused kernel vs.\ exp + add + ln = 3 kernel calls.
  \item \textbf{Existing usage}: many ML libraries implement \texttt{softplus} and
    \texttt{log\_sum\_exp} as fused functions; LEAd is a generalization of softplus
    where the second term is not restricted to $e^y$.
\end{itemize}

\textbf{Verdict}: LEAd has genuine library value. It should be implemented as a
fused function in the monogate library, independent of whether it's a ``primitive.''

\subsection{Criterion 4: Theorem-Statement Coherence}

If LEAd is in F16, the statement ``LSE costs 4n in F16'' becomes false —
a theorem in the current paper stack. Changing F16 to include LEAd would require:
\begin{enumerate}
  \item Revising LSE\_Lower\_Bound.tex (the 4n upper bound is based on the non-LEAd model)
  \item Revising SuperBEST\_Final\_Taxonomy\_Lock.tex (LSE row)
  \item Revising every mention of LSE cost in prior papers
  \item Reconsidering which CONJ\_NO\_OP\_24 proof steps assumed standard F16 convention
\end{enumerate}
CONJ\_NO\_OP\_24\_Closure.tex proves no 24th operator under the PGC+AIT filter,
where PGC includes the standard argument convention. If LEAd is promoted, CONJ\_NO\_OP\_24
would technically be violated — but this is a definitional change, not a mathematical error.

\textbf{Verdict}: promoting LEAd to F16 requires major revision of the theorem stack.
Keeping it separate is simpler and avoids retroactive corrections.

\subsection{Criterion 5: Precedent in the 23-Op Taxonomy}

The 7 extended operators (EEM, EED, EES, EEA, LLA, LLS, LLD) were included because:
\begin{enumerate}
  \item They pass PGC (outside the F16 orbit, not reducible to $\leq 2$ F16 nodes)
  \item They pass AIT (algebraically independent from existing operators on generic inputs)
  \item They arise naturally in exp-log algebra (products/ratios of exp and log terms)
\end{enumerate}

LEAd passes a modified PGC: it is outside the F16 orbit and provides genuine savings.
But it \emph{differs structurally} from all 23 existing operators: its second argument
is not of exp-ln type. It introduces a new \emph{argument-type convention}, not
just a new combination of exp-ln operations.

\textbf{Verdict}: LEAd is categorically different from the 7 extended operators.
Promoting it to the 23-op set would create a 24th operator that breaks CONJ\_NO\_OP\_24.
It belongs in a separate structural-extension category.

\section{Decision}

\begin{theorem}[LEAd architectural decision]
\label{thm:decision}
\textbf{LEAd is not promoted to the standard F16 or 23-op extended primitive sets.}

Rationale:
\begin{enumerate}
  \item LEAd breaks G$_{16}$ group symmetry (Criterion 1) — the defining property of F16.
  \item The theorem stack is clean and consistent without LEAd in F16 (Criterion 4).
  \item CONJ\_NO\_OP\_24 is a proved theorem; adding LEAd to F16 would create a
    definitional violation, requiring the theorem to be re-scoped.
  \item The 23-op extended model already achieves $N_{23}(\mathrm{LSE}) = 2$ via
    EEA+ln (Category C, notation-only); LEAd provides the same cost but via a
    different mechanism.
\end{enumerate}

\textbf{Instead, LEAd is classified as:}
\begin{itemize}
  \item \textbf{Layer 2 structural extension}: available for use in library code
    and implementation guides, labeled as ``LEAd structural extension.''
  \item \textbf{Not counted as 1n in Layer 1}: Layer 1 cost of LEAd = 2n
    (it is exp+add, approximately — the raw-argument structure still requires 2 layer-1 primitives).
  \item \textbf{Genuine numerical benefit}: implement as a fused library function
    (\texttt{monogate.lead(a,b) = log(exp(a)+b)}) for stability and performance.
\end{itemize}
\end{theorem}

\section{Impact Summary}

\begin{center}
\renewcommand{\arraystretch}{1.2}
\begin{tabular}{lccc}
\toprule
Model & LSE cost & Taxonomy & Theorem coherence \\
\midrule
Standard F16 (final) & 4n & 16 operators; G$_{16}$ clean & Full \\
23-op extended & 2n (EEA+ln) & 23 operators; closed & Full (labeled) \\
F16$^+$ (rejected) & 2n (LEAd) & 17 operators; G$_{16}$ broken & Requires revision \\
Library (F16 + LEAd fused) & 2n (LEAd) & Not a formal model & N/A \\
\bottomrule
\end{tabular}
\end{center}

\section{Policy Update}

The following policy additions are appended to the two-layer accounting policy:

\begin{enumerate}
  \item \textbf{LEAd is not in F16}: $\mathrm{LEAd}(a,b) = \ln(e^a+b)$ is outside
    the G$_{16}$ orbit. LSE costs 4n in F16 (Layer 1).
  \item \textbf{LEAd in library}: The monogate library should implement
    \texttt{lead(a,b)} as a fused primitive for numerical stability. Library
    documentation may quote LSE $= 2n$ using LEAd, labeled
    ``using LEAd structural extension.''
  \item \textbf{CONJ\_NO\_OP\_24 scope}: The theorem covers the standard
    argument convention (second argument via $\ln$). LEAd (raw second argument)
    is in a separate category outside the theorem's scope — the theorem is not violated.
  \item \textbf{Future F16$^+$}: If the monogate framework ever adopts a
    ``typed argument'' extension allowing raw-type second arguments, LEAd would
    become a valid primitive. This would create a new model (F16$^+$) with different
    theorems. All F16$^+$ claims must be labeled accordingly.
\end{enumerate}

\end{document}
